% 4.2 具有不变子群的群
\section{具有不变子群的群}

本节考虑具有真不变子群 $\mathbf{T}$ 的群 $\mathbf{G}$ 的若干性质。（若 $\mathbf{T}$ 是 Abel 群，则 $\mathbf{G}$ 具有与空间群相似的结构；例如见式 (3.5.2)。）

我们把 $\mathbf{G}$ 关于 $\mathbf{T}$ 的左陪集展开：
\begin{equation*}
\mathbf{G} = \sum_{x}r_{x}\mathbf{T},
\tag{4.2.1}
\end{equation*}
其中一如既往地假定左陪集代表 $r_{x}$ 一次性选定。我们总取 $r_{1} = e$（$\mathbf{G}$ 的恒等元），并把 $\mathbf{T}$ 的恒等元记作 $t_{1} = e$。

若能把左陪集代表选成一个群 $\mathbf{F}$，我们将假定已作了这样的选择。\footnote{常常有这样的情形：在左陪集代表 $r_{x}$ 的各种选法中，有些选法构成群，有些不构成。例如点群 $432\ (O)$ 可以展开为
\begin{center}
$E\mathbf{T} + C^{+}_{31}\mathbf{T} + C^{-}_{31}\mathbf{T} + C_{2b}\mathbf{T} + C_{2e}\mathbf{T} + C_{2f}\mathbf{T}$
\end{center}
其中 $\mathbf{T}$ 是点群 $222\ (D_{2})\ (= E + C_{2x} + C_{2y} + C_{2z})$；它也可以展开为
\begin{center}
$E\mathbf{T} + C^{+}_{31}\mathbf{T} + C^{-}_{31}\mathbf{T} + C_{2a}\mathbf{T} + C_{2c}\mathbf{T} + C_{2d}\mathbf{T}$。
\end{center}
在第一种情形，$r_{\alpha}$ 构成点群 $32\ (D_{3})$（诚然取向非标准）。（——中译者注：本注按原书脚注照录，第二式中 $C^{+}_{31}$ 的上标疑为 $C^{\pm}_{31}$ 之误排。）}

\noindent\textbf{定义 4.2.1}\quad 若群 $\mathbf{G}$ 关于不变子群 $\mathbf{T}$ 的左陪集代表 $r_{\alpha}$ 可以选成一个子群 $\mathbf{F}$，则称 $\mathbf{G}$ 是 $\mathbf{T}$ 与 $\mathbf{F}$ 的\textit{半直积}（semi-direct product），记作 $\mathbf{G} = \mathbf{T} \wedge \mathbf{F}$。

（半直积的等价定义见定义 1.2.17。）然而，无论 $\mathbf{G}$ 是不是半直积，本节的内容都成立。

把 $\mathbf{G}$ 的元素记作 $(r_{\alpha}t_{a})$（这些记号不要与 Seitz 空间群算符混淆）。其乘法规则为
\begin{equation*}
(r_{\alpha}t_{a})(r_{\beta}t_{b}) = (r_{\alpha\beta}\, t_{\alpha\beta}[t_{a}]_{\beta}t_{b}),
\tag{4.2.2}
\end{equation*}
其中由于 $\mathbf{T}$ 不变，$[t_{a}]_{\beta} = r^{-1}_{\beta}t_{a}r_{\beta} \in \mathbf{T}$，且
\begin{equation*}
r_{\alpha}r_{\beta} = r_{\alpha\beta}t_{\alpha\beta},
\tag{4.2.3}
\end{equation*}
式 (4.2.3) 右端是元素 $r_{\alpha}r_{\beta} \in \mathbf{G}$ 分解为选定的左陪集代表之一 $r_{\alpha\beta}$ 与元素 $t_{\alpha\beta} \in \mathbf{T}$ 之积的唯一分解。对半直积群，对一切 $\alpha, \beta$ 有 $t_{\alpha\beta} = t_{1}$。

$\mathbf{G}$ 的恒等元是 $(r_{1}t_{1})$，由式 (4.2.2) 与 (4.2.3) 得到下列关系：
\begin{equation*}
[t_{a}]_{1} = t_{a}, \quad [t_{1}]_{\alpha} = t_{1}, \quad r_{\alpha 1} = r_{1\alpha} = r_{\alpha}, \quad\text{以及}\quad t_{\alpha 1} = t_{1\alpha} = t_{1};
\tag{4.2.4}
\end{equation*}
\begin{equation*}
[t_{a}]_{\alpha}[t_{b}]_{\alpha} = [t_{a}t_{b}]_{\alpha}
\tag{4.2.5}
\end{equation*}
以及
\begin{equation*}
[t^{-1}_{a}]_{\alpha} = \left([t_{a}]_{\alpha}\right)^{-1}.
\tag{4.2.6}
\end{equation*}
于是由群 $\mathbf{G}$ 的结合性（利用式 (4.2.2) 与 (4.2.5)）推得，对一切 $\alpha, \beta, \gamma$
\begin{equation*}
r_{\alpha\beta, \gamma} = r_{\alpha, \beta\gamma}
\tag{4.2.7}
\end{equation*}
以及
\begin{equation*}
t_{\alpha\beta, \gamma}[t_{\alpha\beta}]_{\gamma}[[t_{a}]_{\beta}]_{\gamma} = t_{\alpha, \beta\gamma}[t_{a}]_{\beta\gamma}t_{\beta\gamma}.
\tag{4.2.8}
\end{equation*}
注意由上述各方程组可得
\begin{equation*}
r_{\alpha, \beta 1} = r_{\alpha\beta, 1} = r_{\alpha\beta} = r_{1, \alpha\beta} = r_{1\alpha, \beta} = r_{\alpha, 1\beta} = r_{\alpha 1, \beta}
\tag{4.2.9}
\end{equation*}
但是，尽管
\begin{equation*}
t_{\alpha, \beta 1} = t_{1\alpha, \beta} = t_{\alpha, 1\beta} = t_{\alpha 1, \beta} = t_{\alpha\beta}, \quad t_{\alpha\beta, 1} = t_{1, \alpha\beta} = t_{1}.
\tag{4.2.10}
\end{equation*}
在式 (4.2.8) 中令 $a = 1$，得到一个对一切 $\alpha, \beta, \gamma$ 都重要的关系：
\begin{equation*}
t_{\alpha\beta, \gamma}[t_{\alpha\beta}]_{\gamma} = t_{\alpha, \beta\gamma}t_{\beta\gamma}.
\tag{4.2.11}
\end{equation*}
以下几点值得注意。$(r_{1}t_{a})$ 的逆是 $(r_{1}t^{-1}_{a})$。给定 $r_{\alpha}$，存在唯一的陪集代表 $r_{\bar{\alpha}}$ 使 $r_{\alpha\bar{\alpha}} = r_{\bar{\alpha}\alpha} = r_{1}$；$(r_{\alpha}t_{1})$ 的逆是 $(r_{\bar{\alpha}}t_{\alpha\bar{\alpha}})$。并且
\begin{equation*}
t_{\bar{\alpha}\alpha} = [t_{\alpha\bar{\alpha}}]_{\alpha},
\tag{4.2.12}
\end{equation*}
以及
\begin{equation*}
\left[[t_{a}]_{\alpha}\right]_{\bar{\alpha}} = t^{-1}_{\alpha\bar{\alpha}}t_{a}t_{\alpha\bar{\alpha}}.
\tag{4.2.13}
\end{equation*}
最后，容易验证
\begin{equation*}
r_{\alpha}t_{a}r^{-1}_{\alpha} = t_{\alpha\bar{\alpha}}[t_{a}]_{\bar{\alpha}}t^{-1}_{\alpha\bar{\alpha}}.
\tag{4.2.14}
\end{equation*}
我们把式 (4.2.14) 的右端取作 $[t_{a}]_{\alpha}^{-1}$ 的定义；它没有别的定义，因为 $r^{-1}_{\alpha}$ 不一定是选定的陪集代表之一。

\noindent\textbf{定理 4.2.1}\quad $(r_{\alpha}t_{a})$ 的逆是 $\left(r_{\bar{\alpha}}[t^{-1}_{a}]_{\bar{\alpha}}t^{-1}_{\alpha\bar{\alpha}}\right)$。

因为
\begin{equation*}
(r_{\alpha}t_{a})(r_{\bar{\alpha}}[t^{-1}_{a}]_{\bar{\alpha}}t^{-1}_{\alpha\bar{\alpha}}) = \left(r_{\alpha\bar{\alpha}}\, t_{\alpha\bar{\alpha}}[t_{a}]_{\bar{\alpha}}[t^{-1}_{a}]_{\bar{\alpha}}t^{-1}_{\alpha\bar{\alpha}}\right),
\end{equation*}
由式 (4.2.5) 及 $r_{\alpha\bar{\alpha}} = r_{1}$ 的事实，它等于恒等元。留给读者一个练习：利用式 (4.2.12) 与 (4.2.13) 验证 $\left(r_{\bar{\alpha}}[t^{-1}_{a}]_{\bar{\alpha}}t^{-1}_{\alpha\bar{\alpha}}\right) = (r_{\alpha}t_{a})^{-1}$ 亦可从 $(r_{\alpha}t_{a})(r_{\bar{\alpha}}t^{-1}_{a}t_{\alpha\bar{\alpha}}) = (r_{1}t_{1})$ 得出。

\noindent\textbf{定理 4.2.2}\quad 陪集代表的下标构成一个 $|\mathbf{G}|/|\mathbf{T}|$ 阶的群，记作 $\mathbf{R}$。它当然同构于因子群 $\mathbf{G}/\mathbf{T}$。

$\mathbf{R}$ 中的乘法规则是：若 $r_{\alpha}r_{\beta} \in r_{\gamma}\mathbf{T}$，则 $\alpha\beta = \gamma$。封闭性由构造即得，结合性由式 (4.2.7) 得到，恒等元是 1（见式 (4.2.4)），$\alpha$ 的逆是 $\bar{\alpha}$。

现在把 4.1 节的结果具体化到子群不变的情形。为此把 4.1 节的 $\mathbf{K}_{1}$、$p_{\alpha}$ 与 $\Gamma^{j}$ 分别认同为本节的 $\mathbf{T}$、$r_{\alpha}$ 与 $\mathbf{D}^{i}$。

$\mathbf{D}^{i}$ 是 $\mathbf{T}$ 的维数为 $d_{i}$ 的表示，基为 $\langle \phi_{r} |$，使得（见式 (4.1.2)）
\begin{equation*}
t_{a}\phi_{r} = \sum_{p}\phi_{p}\mathbf{D}^{i}(t_{a})_{pr}.
\tag{4.2.15}
\end{equation*}
函数 $\phi_{\alpha p} = r_{\alpha}\phi_{p}$（$\alpha = 1$ 到 $|\mathbf{G}|/|\mathbf{T}|$，$p = 1$ 到 $d_{i}$）张成一个在 $\mathbf{G}$ 下不变的矢量空间 $\mathbf{Q}$（见定理 4.1.2）。子群 $\mathbf{T}_{\alpha} = \{r_{\alpha}t_{a}r^{-1}_{\alpha};\ t_{a} \in \mathbf{T}\}$ 在 $\Omega_{\alpha}$ 上由
\begin{equation*}
t_{a}\phi_{\alpha r} = t_{a}r_{\alpha}\phi_{r} = r_{\alpha}[t_{a}]_{\alpha}\phi_{r} = \sum_{p}\phi_{\alpha p}\mathbf{D}^{i}([t_{a}]_{\alpha})_{pr}
\tag{4.2.16}
\end{equation*}
表示。式 (4.2.16) 蕴含
\begin{equation*}
\mathbf{D}^{i}_{\alpha}(t_{a}) = \mathbf{D}^{i}([t_{a}]_{\alpha}).
\tag{4.2.17}
\end{equation*}
由于此式对一切 $t_{a} \in \mathbf{T}$ 成立，称 $\mathbf{D}^{i}_{\alpha}$ 在 $\mathbf{G}$ 下\textit{共轭}于 $\mathbf{D}^{i}$。对 $\alpha = 1$ 到 $|\mathbf{G}|/|\mathbf{T}|$ 的表示集合 $\{D^{i}_{\alpha}\}$ 构成一族共轭表示。两个共轭表示可能等价；不过这并不影响随后的论证。

从一族共轭表示中选出的一组彼此不等价的表示的极大集合称为一个\textit{星}（star）（有些作者称之为轨道（orbit））。很容易看出这与已对空间群给出的定义一致。若把 $t_{a}$ 认同为平移 $\{E \mid \mathbf{t}\}$、把 $r_{\alpha}$ 认同为操作 $\{R \mid \mathbf{v}\}$，则
\begin{equation*}
[t_{a}]_{\alpha} \rightarrow \{R \mid \mathbf{v}\}^{-1}\{E \mid \mathbf{t}\}\{R \mid \mathbf{v}\} = \{E \mid R^{-1}\mathbf{t}\}.
\tag{4.2.18}
\end{equation*}
若把 $\mathbf{D}^{i}$ 认同为 $\Delta^{\mathbf{k}}$，则由式 (4.2.17)、(4.2.18)、(3.4.3) 与 (3.6.2)
\begin{equation*}
\begin{aligned}
\mathbf{D}^{i}_{\alpha}(t_{a}) &= \mathbf{D}^{i}([t_{a}]_{\alpha}) \rightarrow \Delta^{\mathbf{k}}(\{E \mid R^{-1}\mathbf{t}\})\\
&= \exp(-\mathrm{i}\mathbf{k} \cdot R^{-1}\mathbf{t})\mathbf{1} = \exp(-\mathrm{i}R\mathbf{k} \cdot \mathbf{t})\mathbf{1}\\
&= \Delta^{R\mathbf{k}}(\{E \mid \mathbf{t}\}) = \mathbf{D}^{i}_{\alpha}(t_{a}).
\end{aligned}
\tag{4.2.19}
\end{equation*}
两种定义的等价性现在来自如下事实：若 $R\mathbf{k} = \mathbf{k}$ 则 $\Delta^{R\mathbf{k}} \equiv \Delta^{\mathbf{k}}$，因此从平移群 $\mathbf{T}$ 的一族共轭表示中选出一组彼此不等价的极大表示集，就是构成一颗星——不管我们把星看作由表示组成还是由 $\mathbf{k}$ 矢量组成。

使 $D^{i}_{\gamma} \equiv D^{i}$ 的那些下标 $\gamma \in \mathbf{R}$ 的集合构成一个群 $\overline{\mathbf{K}}^{i}$，称为 $D^{i}$ 在 $\mathbf{G}$ 中的小陪群。与空间群术语（见定义 3.7.3）的对应是直接的：那里 $\mathbf{k}$ 在 $\mathbf{G}$ 中的小陪群 $\overline{\mathbf{G}}^{\mathbf{k}}$ 满足
\begin{equation*}
\begin{aligned}
\mathbf{D}^{i}_{\alpha\bar{\gamma}\bar{\alpha}}(t_{a}) &= \mathbf{D}^{i}([t_{a}]_{\alpha\bar{\gamma}\bar{\alpha}}) = \mathbf{D}^{i}\left(\left[[t_{a}]_{\alpha\bar{\gamma}\bar{\alpha}}\right]_{\alpha}\right)\\
&= \mathbf{D}^{i}([t_{a}]_{\alpha\gamma}) = \mathbf{D}^{i}_{\gamma}([t_{a}]_{\alpha}) = \mathbf{D}^{i}_{\gamma\alpha}(t_{a}),
\end{aligned}
\tag{4.2.19a}
\end{equation*}
在空间群理论中相应的结果是 $\overline{\mathbf{G}}^{R\mathbf{k}} = \{RSR^{-1};\ S \in \overline{\mathbf{G}}^{\mathbf{k}}\}$。

若一族共轭表示的所有成员互不等价，则小陪群仅由 1 组成。这对应于 $\mathbf{k}$ 是布里渊区的一般点。由方程
\begin{equation*}
\mathbf{K}^{i} = \sum_{\gamma \in \overline{\mathbf{K}}^{i}}r_{\gamma}\mathbf{T}
\tag{4.2.20}
\end{equation*}
定义的群 $\mathbf{K}^{i}$（其中 $\gamma$ 取遍小陪群）称为 $D^{i}$（$i$）在 $\mathbf{G}$ 中的\textit{小群}。小陪群同构于因子群 $\mathbf{K}^{i}/\mathbf{T}$。

现在考察由空间 $\mathbf{Q}$ 确定的诱导表示 $\mathbf{D}^{i} \uparrow \mathbf{G}$。与式 (4.1.4) 相应的方程是
\begin{equation*}
(r_{\alpha}t_{a})\phi_{\tau r} = \sum_{p}\phi_{\sigma p}\mathbf{D}^{i}(t_{\alpha\tau}[t_{a}]_{\tau})_{pr}
\tag{4.2.21}
\end{equation*}
其中 $\sigma = \alpha\tau$。

与式 (4.1.9) 相应的方程是
\begin{equation*}
\chi(r_{\alpha}t_{a}) = \delta_{\alpha 1}\sum_{\tau}\chi^{i}_{\tau}(t_{a}).
\tag{4.2.22}
\end{equation*}
也就是说，在 $\mathbf{D}^{i} \uparrow \mathbf{G}$ 中只有 $\mathbf{T}$ 的元素才有非零特征标；而在 $\mathbf{T}$ 中，每个元素的特征标是它在与 $D^{i}$ 共轭的一组完全表示上取的特征标之和。

由于 $\mathbf{T}$ 不变，Johnston 不可约性判据（定理 4.1.4）的推论 (i) 适用，并断言：若与 $D^{i}$ 共轭的表示集合中没有两个等价，则 $\mathbf{D}^{i} \uparrow \mathbf{G}$ 不可约。于是对空间群 $\mathbf{G}$，若 $\mathbf{k}$ 是布里渊区的一般点，则 $\Delta^{\mathbf{k}} \uparrow \mathbf{G}$ 不可约。

由式 (4.2.22) 可见
\begin{equation*}
(\mathbf{D}^{i} \uparrow \mathbf{G}) \downarrow \mathbf{T} = \sum_{\tau}D^{i}_{\tau}.
\tag{4.2.23}
\end{equation*}
当 $\mathbf{D}^{i} \uparrow \mathbf{G}$ 不可约时，$D^{i}_{\tau}$（$\tau = 1$ 到 $|\mathbf{G}|/|\mathbf{T}|$）彼此不等价，式 (4.2.23) 右端等于含 $D^{i}$ 的那颗星（简称星 $i$）。并且由 Frobenius 互反定理 4.1.5 得：$D^{i}_{\tau} \uparrow \mathbf{G}$ 与 $\mathbf{D}^{i} \uparrow \mathbf{G}$ 等价，因为按该定理 $D^{i}_{\tau} \uparrow \mathbf{G}$ 中含 $\Gamma$（任意一个含于 $\mathbf{D}^{i} \uparrow \mathbf{G}$ 的不可约表示）的次数等于 $\Gamma$ 限制后含 $D^{i}_{\tau}$ 的次数，与 $i$ 无关。

一般情形，当 $\mathbf{D}^{i} \uparrow \mathbf{G}$ 可约时，小陪群由 1 以外的元素组成。由式 (4.2.20) 可见它由 $|\mathbf{K}^{i}|/|\mathbf{T}|$ 个元素组成；这意味着在和 $\sum_{\tau}D^{i}_{\tau}$ 中，星 $i$ 将出现 $|\mathbf{K}^{i}|/|\mathbf{T}|$ 次。

至此我们已证明：当小陪群只含一个元素时，表示 $\mathbf{D}^{i} \uparrow \mathbf{G}$ 不可约，而且它不仅可以由 $D^{i}$ 诱导，也可以由星 $i$ 的任一成员诱导。这样的不可约表示限制到 $\mathbf{T}$ 上时，星中的每个成员恰好出现一次。并非 $\mathbf{G}$ 的每个表示都能这样得到，因为平移群存在含等价表示的共轭表示集合。这后一种情形正是接下来要讨论的。
